% ===================================================================
%  TAVOLA N. 6 — La casa di dentro. --- Mobili. --- Cibo. --- Luce.
%  Traduction 2026 d'apres texto/io, controlee sur texto/fr et
%  texto/en. Consigne : voir texto/it/10-tavola-01.tex.
%
%  Deux notes marquees toutes deux « (*) », et deux appels : celui de
%  « babo » au bloc c1-12-1, celui de « lambrequino » au bloc c2-01-1.
%  L'ordre les departage. La femme de chambre du bloc c1-06-1 n'a pas
%  de gros plan : l'ido porte « (6) », qui est l'eponge, le francais
%  « (16) », qui a raison, et la planche ne porte pas de 16 releve.
%
%  DEUX MEUBLES DE CETTE PAGE PORTENT LE NOM D'UN DANGER QUI N'EXISTE
%  PLUS.
%
%  1. LA CREDENZA DE L'ALINEA 2 DE LA TROISIEME COLONNE (5). Le
%     buffet de la salle a manger s'appelle en italien
%     credenza, et « credenza » est LA CROYANCE, la
%     confiance — le meme mot que « credere ». C'etait la table sur
%     laquelle un serviteur GOUTAIT les plats devant le maitre pour
%     prouver qu'ils n'etaient pas empoisonnes ; « far la credenza »
%     etait le nom de ce geste. Le meuble a garde le nom du rite.
%     L'allemand dit « Anrichte », le lieu ou l'on prepare ;
%     l'anglais « sideboard », par sa place dans la piece ; le
%     francais « buffet », dont personne ne sait l'origine. L'italien
%     seul nomme ce meuble par une peur.
%
%  2. L'ARMADIO DE L'ALINEA 5 (14). Il vient de « armarium », qui
%     etait chez les Romains L'ARMOIRE AUX ARMES. Le francais l'a pris
%     aussi — armoire — mais avec un O qui a brouille la piste ;
%     l'italien a garde l'A, et le mot dit encore ce qu'il disait.
%
%  LE TABLEAU 6 ALLEMAND AVAIT TROUVE QUE « MÖBEL » ETAIT FRANCAIS ET
%  QUE PAS UN MEUBLE NE L'ETAIT — l'abstraction empruntee, les membres
%  indigenes. Ici la question ne se pose meme pas : « mobili » est le
%  latin « mobilia », les choses qui se deplacent, et c'est le mot
%  d'origine, celui que le francais a emprunte et que l'allemand a
%  emprunte au francais. LA COLONNE ITALIENNE EST TOUJOURS EN AMONT.
%
%  ET LA NOTE (*) DU LAMBREQUIN OUBLIE L'ITALIEN, comme l'autre note
%  oubliait l'allemand. Guignon aligne le francais, l'espagnol, le
%  portugais, l'allemand et l'anglais, qui disent tous
%  « lambrequin ». L'italien dit mantovana — DE MANTOUE, la
%  ville. C'est la seconde fois que cette colonne rencontre un objet
%  domestique nomme par le lieu d'ou il venait : le tableau 1 avait la
%  lavagna, l'ardoise du village de Lavagna. Deux fois sur
%  six tableaux, et jamais dans les trente-trois autres colonnes.
%
%     UNE LANGUE PARLEE PAR DES MARCHANDS DE VILLE EN VILLE NOMME LES
%     CHOSES PAR LA VILLE D'OU ELLES ARRIVENT. C'est la regle du
%     tableau 11 allemand — la langue d'une institution est celle du
%     lieu qui l'a batie — mais pour des objets, et a l'echelle non
%     d'un pays mais d'un bourg.
% ===================================================================

%%K t06-apar-1 apar
\VUsaut{6.0mm}
\VUtitre{15.0pt}{60}{TAVOLA N. 6}
\VUsaut{1.4mm}
\VUtitre{11.4pt}{0}{La casa di dentro, i mobili, il cibo,}
\VUsaut{0.4mm}
\VUtitre{11.4pt}{0}{la luce.}
\VUsaut{2.2mm}

%%K t06-apar-2 apar
La casa è finita. Lo zio di Giovanni si è trasferito nella sua nuova abitazione, di
cui conosciamo già la distribuzione. Vediamo adesso come l'ha ammobiliata. Per prima
cosa visitiamo:

%%K t06-tit-1 sub
\VUpk{10.2pt}{40}{La camera da letto.}
\VUsaut{3.0mm}

%%K t06-c1-01-1 p
1. --- Siamo arrivati in un brutto momento: la domestica sta pulendo la
\VUgras{stanza}.

%%K t06-c1-02-1 p
2. --- Sul \VUgras{lavabo} \textsuperscript{(1)}, sparse qua e là, ci sono le cose con
cui ci si lava, ci si pettina e ci si prepara. Sono il \VUgras{catino}
\textsuperscript{(2)} e la \VUgras{brocca} \textsuperscript{(3)}, la
\VUgras{spazzola} \textsuperscript{(4)}, il \VUgras{pettine} \textsuperscript{(5)}, il
\VUgras{sapone} \textsuperscript{(6)} e la \VUgras{spugna} \textsuperscript{(7)}, e
infine una \VUgras{boccetta} \textsuperscript{(8)} di collutorio.

%%K t06-c1-03-1 p
3. --- Per terra, davanti al lavabo, c'è il \VUgras{secchio} \textsuperscript{(9)} per
l'acqua sporca.

%%K t06-c1-04-1 p
4. --- Nell'\VUgras{anticamera} \textsuperscript{(10)} vediamo qualche
\VUgras{attaccapanni} \textsuperscript{(13)} e due \VUgras{ombrelli}
\textsuperscript{(11)} nel \VUgras{portaombrelli} \textsuperscript{(12)}.

%%K t06-c1-05-1 p
5. --- Dall'altra parte della porta c'è l'\VUgras{armadio a specchio}
\textsuperscript{(14)}, dove sta la \VUgras{biancheria} \textsuperscript{(15)}.

%%K t06-c1-06-1 p
6. --- Mentre la \VUgras{cameriera} \textsuperscript{(16)} rifà il \VUgras{letto}
\textsuperscript{(17)}, guardiamone le varie parti. La \VUgras{lettiera}
\textsuperscript{(20)} porta un buon \VUgras{materasso a molle}
\textsuperscript{(21)}, sul quale Giustina metterà fra poco un \VUgras{materasso} di
lana \textsuperscript{(22)}. Poi prenderà dalle sedie i due \VUgras{lenzuoli}
\textsuperscript{(23)} e la \VUgras{coperta} \textsuperscript{(24)}. Poi metterà a capo
del letto il \VUgras{capezzale} \textsuperscript{(25)} e il \VUgras{guanciale}
\textsuperscript{(26)}, che stanno con il \VUgras{piumino} \textsuperscript{(27)} sullo
\VUgras{scendiletto} \textsuperscript{(32)}. Passa il piumino da polvere sulle
\VUgras{tende} \textsuperscript{(19)} e sul \VUgras{baldacchino}
\textsuperscript{(18)}, e una parte del lavoro è fatta.

%%K t06-c1-07-1 p
7. --- Poi dovrà mettere un po' in ordine il \VUgras{comodino}
\textsuperscript{(28)}, caricare la \VUgras{sveglia} \textsuperscript{(29)}, preparare
il \VUgras{lumino da notte} \textsuperscript{(30)} e portare via la \VUgras{candela}
\textsuperscript{(31)} per pulire il candeliere.

%%K t06-tit-2 sub
\VUsaut{5.0mm}
\VUpk{10.2pt}{40}{Il cestino da lavoro e il camino.}
\VUsaut{3.0mm}

%%K t06-c1-08-1 p
8. --- Anche il \VUgras{cestino da lavoro} \textsuperscript{(34)} accanto allo
\VUgras{sgabello} \textsuperscript{(33)} ha un gran bisogno di essere riordinato.
Dovrà piegare il \VUgras{ricamo} \textsuperscript{(35)}, riporre nel
\VUgras{tavolino da lavoro} \textsuperscript{(38)} il \VUgras{ditale}, il
\VUgras{filo} \textsuperscript{(36)}, gli \VUgras{aghi} \textsuperscript{(37)} e le
\VUgras{forbici} \textsuperscript{(39)}, e chiudere il coperchio del tavolino con la
\VUgras{chiave} \textsuperscript{(40)}.

%%K t06-c1-09-1 p
9. --- Sul \VUgras{tavolino rotondo} \textsuperscript{(41)} in mezzo alla stanza
Giustina cambierà l'acqua della \VUgras{caraffa} \textsuperscript{(42)}. Metterà lo
\VUgras{zucchero} nella \VUgras{zuccheriera} \textsuperscript{(43)}, pulirà le
\VUgras{mollette} \textsuperscript{(44)} e porterà via la \VUgras{spazzola per abiti}
\textsuperscript{(46)}.

%%K t06-c1-10-1 p
10. --- La parte destra della stanza le darà meno lavoro, perché la
\VUgras{chaise-longue} \textsuperscript{(47)} e il suo \VUgras{cuscino}
\textsuperscript{(48)} sono al loro posto, e siccome non fa freddo non è stato spostato
nulla di quello che appartiene al \VUgras{camino} \textsuperscript{(49)}. La
\VUgras{paletta} \textsuperscript{(50)}, le \VUgras{molle} \textsuperscript{(51)}, gli
\VUgras{alari} \textsuperscript{(55)} e il \VUgras{parascintille}
\textsuperscript{(56)} sono puliti; il \VUgras{paravento del camino}
\textsuperscript{(54)} non è stato spostato e la \VUgras{cassa della legna}
\textsuperscript{(52)} è ben fornita di \VUgras{legna da ardere}
\textsuperscript{(53)}.

%%K t06-noto-1 noto
\VUnotes{143.2mm}{%
(*) Babo = un bambino piccolo; inglese \textit{baby}, francese \textit{bébé},
italiano \textit{bambino}.%
}

%%K t06-noto-2 noto
\VUnotes{143.2mm}{%
(*) Lambrequino = francese \textit{lambrequin}, spagnolo \textit{lambrequines},
portoghese \textit{lambrequins}, e anche tedesco e inglese \textit{lambrequins} ---
dunque tedesco, inglese, francese, portoghese e spagnolo allo stesso modo.%
}

%%K t06-c1-11-1 p
11. --- Dovrà ancora spolverare l'\VUgras{orologio da camino} \textsuperscript{(57)}, i
\VUgras{candelabri} \textsuperscript{(58)} e i \VUgras{piattini da ninnoli}
\textsuperscript{(59)}, e poi pulire lo \VUgras{specchio} \textsuperscript{(60)}.

%%K t06-c1-12-1 p
12. --- Quanto alla \VUgras{culla} \textsuperscript{(61)} del bambino (il babo (*)), è
già pronta.

%%K t06-tit-3 sub
\VUsaut{5.0mm}
\VUpk{10.2pt}{40}{Il salotto.}
\VUsaut{3.0mm}

%%K t06-c2-01-1 p
1. --- Passiamo adesso al salotto. È una \VUgras{stanza} molto bella. Il camino,
nell'angolo di destra, è ornato di una delicata \VUgras{statuetta}
\textsuperscript{(1)} e di due bei \VUgras{vasi} \textsuperscript{(3)}, e drappeggiato
con una \VUgras{mantovana} \textsuperscript{(2)} di velluto (*).

%%K t06-c2-02-1 p
2. --- Perché le scintille del focolare non diano fuoco al tappeto, davanti è stato
messo un \VUgras{parafuoco} \textsuperscript{(4)} di rame.

%%K t06-c2-03-1 p
3. --- Accanto, aperto, c'è un elegante \VUgras{scrittoio da signora}
\textsuperscript{(5)}. È lì che la \VUgras{padrona di casa} \textsuperscript{(23)}
scrive le sue lettere.

%%K t06-c2-04-1 p
4. --- La finestra è vestita di \VUgras{tendine} \textsuperscript{(6)} tenute indietro
da \VUgras{fermatende} \textsuperscript{(8)} agganciati ai \VUgras{ganci}
\textsuperscript{(7)}, e di pesanti tende di stoffa sormontate da una
\VUgras{balza} \textsuperscript{(9)}.

%%K t06-c2-05-1 p
5. --- Nel vano della finestra una \VUgras{fioriera} \textsuperscript{(11)} regge una
\VUgras{pianta} \textsuperscript{(10)} verde, una magnifica palma le cui foglie
arrivano quasi alla \VUgras{lampada a petrolio} \textsuperscript{(12)}.

%%K t06-c2-06-1 p
6. --- Il \VUgras{paralume} \textsuperscript{(13)} l'ha fatto Maria, l'incantevole
\VUgras{signorina} \textsuperscript{(14)} seduta al \VUgras{pianoforte}
\textsuperscript{(15)}. Sta molto dritta sullo \VUgras{sgabello}
\textsuperscript{(16)} e ripassa un pezzo. È molto brava e molto ordinata, come mostra
il suo \VUgras{mobile della musica} \textsuperscript{(17)}: è in ordine perfetto.

%%K t06-tit-4 sub
\VUsaut{5.0mm}
\VUpk{10.2pt}{40}{Il monello.}
\VUsaut{3.0mm}

%%K t06-c2-07-1 p
7. --- Suo fratello Paolo cerca un \VUgras{libro} \textsuperscript{(19)} nella
\VUgras{libreria} \textsuperscript{(18)}. È un \VUgras{ragazzo}
\textsuperscript{(20)} irrequieto e proprio insopportabile. Si aggrappa alla libreria
per arrivare al libro che sta troppo in alto, e rischia di far cadere il
\VUgras{busto} \textsuperscript{(21)} che c'è sopra.

%%K t06-c2-08-1 p
8. --- Riesce a farlo solo perché sua madre sta parlando con un'amica, una
\VUgras{signora} \textsuperscript{(22)} venuta a passare qualche giorno da lei. La
madre è seduta sul \VUgras{divano} \textsuperscript{(24)} accanto alla
\VUgras{poltrona} \textsuperscript{(25)}, e la sua amica sulla \VUgras{sedia}
\textsuperscript{(27)}, di cui vediamo solo lo \VUgras{schienale}
\textsuperscript{(26)}.

%%K t06-c2-09-1 p
9. --- La grande \VUgras{cornice} \textsuperscript{(30)} appesa al muro contiene il
\VUgras{ritratto} \textsuperscript{(29)} di un parente. Nell'\VUgras{album}
\textsuperscript{(32)} posato sul tavolo sono raccolte le \VUgras{fotografie}
\textsuperscript{(33)} del resto della famiglia. I bambini lo sfogliavano poco fa: le
\VUgras{sedie} \textsuperscript{(31)} sono ancora raccolte intorno al tavolo.

%%K t06-c2-10-1 p
10. --- Il \VUgras{vaso cinese} \textsuperscript{(34)} di porcellana accanto al
\VUgras{tagliacarte} \textsuperscript{(35)} è stato regalato a Maria per il suo
onomastico, e il \VUgras{paravento} \textsuperscript{(36)} che riempie l'angolo della
stanza l'ha portato dalla Cina suo zio.

%%K t06-c2-11-1 p
11. --- Rallegrato dai bei \VUgras{quadri} \textsuperscript{(37)} appesi alla
\VUgras{parete} \textsuperscript{(42)}, illuminato dal \VUgras{lampadario}
\textsuperscript{(38)} del soffitto, le cui \VUgras{lampadine elettriche}
\textsuperscript{(39)} la sera brillano così allegre, questo salotto ha un'aria molto
gradevole. Il morbido \VUgras{tappeto} \textsuperscript{(40)} steso sul pavimento e il
\VUgras{campanello elettrico} \textsuperscript{(41)}, così comodo, lo rendono
perfettamente confortevole.

%%K t06-tit-5 sub
\VUsaut{3.6mm}
\VUpk{10.2pt}{40}{La sala da pranzo.}
\VUsaut{3.0mm}

%%K t06-c3-01-1 p
1. --- Hanno suonato per il pranzo. Tutta la famiglia, tranne Maria, è riunita intorno
alla tavola. La sala da pranzo è piacevolmente scaldata da un'ottima \VUgras{stufa}
\textsuperscript{(1)} di maiolica con un \VUgras{tubo} \textsuperscript{(2)} sottile.

%%K t06-c3-02-1 p
2. --- In questa stanza non ci sono molti mobili. Lungo la parete, sopra lo
\VUgras{sgabello} \textsuperscript{(4)}, è appesa una graziosa \VUgras{fontanella}
\textsuperscript{(3)}. Lì accanto c'è la \VUgras{credenza} \textsuperscript{(5)}, su
cui si posano i piatti freddi, i dolci e le varie stoviglie che al momento non
servono.

%%K t06-c3-03-1 p
3. --- Così sul ripiano di sopra vediamo un \VUgras{pollo arrosto}
\textsuperscript{(7)}, un \VUgras{pesce} \textsuperscript{(8)} e una \VUgras{coppa}
\textsuperscript{(10)} di \VUgras{frutta} \textsuperscript{(9)}. Sotto ci sono il
\VUgras{formaggio} \textsuperscript{(11)}, il \VUgras{pane} \textsuperscript{(12)} e
l'\VUgras{oliera} \textsuperscript{(13)}, che contiene tutto quello che serve per
condire l'insalata: l'olio, l'aceto, il \VUgras{sale} \textsuperscript{(14)} e il
\VUgras{pepe} \textsuperscript{(15)}. Infine, sul ripiano più basso, c'è una
\VUgras{torta} \textsuperscript{(17)} accanto al \VUgras{vasetto della senape}
\textsuperscript{(16)}.

%%K t06-c3-04-1 p
4. --- L'orologio della sala da pranzo, che si chiama \VUgras{orologio a muro}
\textsuperscript{(20)}, ha una forma molto insolita. È incassato in una cornice di
legno che porta anche un \VUgras{barometro} \textsuperscript{(19)} e un
\VUgras{termometro} \textsuperscript{(18)}.

%%K t06-tit-6 sub
\VUsaut{4.6mm}
\VUpk{10.2pt}{40}{Il pranzo viene servito.}
\VUsaut{3.0mm}

%%K t06-c3-05-1 p
5. --- Tutt'intorno alla stanza le pareti sono rivestite di \VUgras{boiserie}
\textsuperscript{(21)} di rovere incerato. Una \VUgras{portiera}
\textsuperscript{(22)} di stoffa chiude abbastanza bene la porta che dà sulla veranda.
È lì che la famiglia andrà a prendere il caffè dopo il pranzo. Le \VUgras{tazze}
\textsuperscript{(24)} e i \VUgras{piattini} \textsuperscript{(25)} sono già pronti sul
\VUgras{mobile delle stoviglie} \textsuperscript{(23)}.

%%K t06-c3-06-1 p
6. --- Sopra la tavola una \VUgras{lampada a sospensione} \textsuperscript{(26)} è
fissata al soffitto; la sera la sua luce vivissima riempie tutta la sala da pranzo.

%%K t06-c3-07-1 p
7. --- Guardiamo adesso la \VUgras{tavola} \textsuperscript{(27)} stessa. Quando la
famiglia è entrata, la tavola era apparecchiata. Sulla \VUgras{tovaglia}
\textsuperscript{(28)} erano stati messi, al posto di ciascuno, un \VUgras{piatto}
\textsuperscript{(33)}, un \VUgras{tovagliolo} \textsuperscript{(29)}, un
\VUgras{cucchiaio} \textsuperscript{(34)}, una \VUgras{forchetta}
\textsuperscript{(35)}, un \VUgras{coltello} \textsuperscript{(36)} e un
\VUgras{bicchiere} \textsuperscript{(32)}. C'erano anche \VUgras{bottiglie}
\textsuperscript{(30)} per il vino e una \VUgras{caraffa} \textsuperscript{(31)} per
l'acqua.

%%K t06-c3-08-1 p
8. --- Il \VUgras{domestico} \textsuperscript{(39)} ha portato prima la
\VUgras{zuppiera} \textsuperscript{(37)}, in cui fuma ancora un po' di minestra.
Adesso serve la prima \VUgras{portata} \textsuperscript{(38)}, un piatto di carne. Poi
farà girare quelle che abbiamo già nominato; e infine, dopo il \VUgras{dolce}, quando
i padroni saranno passati in veranda, sparecchierà la tavola e rimetterà in ordine la
sala da pranzo.

%%K t06-c3-08-2 p
\textit{Nota.} --- \textit{Le parole comuni che riguardano il cibo sono indicate qui
solo per sommi capi. L'elenco sarà completato nelle due tavole «L'albergo» (n. 13) e
«Il mercato» (n. 15).}

%%K t06-tit-7 sub
\VUsaut{4.6mm}
\VUpk{10.2pt}{40}{La cucina.}
\VUsaut{3.0mm}

%%K t06-c4-01-1 p
1. --- Nella cucina, che intravediamo dalla porta socchiusa, c'è un disordine
pittoresco. Davanti al \VUgras{focolare} \textsuperscript{(1)}, dove crepita il
\VUgras{fuoco} \textsuperscript{(3)}, è sistemato il \VUgras{girarrosto}
\textsuperscript{(5)}. Un bell'\VUgras{arrosto} \textsuperscript{(7)} è sullo
\VUgras{spiedo} \textsuperscript{(6)} e sta finendo di cuocere.

%%K t06-c4-02-1 p
2. --- Il \VUgras{soffietto} \textsuperscript{(8)} e la \VUgras{paletta del carbone}
\textsuperscript{(9)}, che sono stati appena usati, sono rimasti per terra insieme ai
\VUgras{ceppi} \textsuperscript{(10)}.

%%K t06-c4-03-1 p
3. --- La mensola del camino è in ordine: la \VUgras{lanterna}
\textsuperscript{(11)}, il \VUgras{portafiammiferi} \textsuperscript{(13)} con i
\VUgras{fiammiferi} \textsuperscript{(12)} dentro, il \VUgras{candeliere}
\textsuperscript{(14)} e la \VUgras{bugia} \textsuperscript{(15)} con la sua
\VUgras{candela} \textsuperscript{(16)} sono tutti al loro posto. Ma il grande
\VUgras{secchio del carbone} \textsuperscript{(18)}, pieno del \VUgras{carbone}
\textsuperscript{(17)} che la \VUgras{cuoca} \textsuperscript{(23)} è appena andata a
prendere in cantina, è rimasto in mezzo alla cucina.

%%K t06-c4-04-1 p
4. --- La verità è che la poveretta deve sbrigarsi se vuole che il pranzo sia pronto
in tempo. Adesso è alla \VUgras{cucina economica} \textsuperscript{(19)} e mescola una
\VUgras{salsa} \textsuperscript{(24)} che non deve attaccarsi.

%%K t06-c4-05-1 p
5. --- Nel \VUgras{forno} \textsuperscript{(20)} cuoce una torta che ha fatto per la
merenda dei bambini. Sopra la cucina economica sono appesi il \VUgras{mestolo}
\textsuperscript{(21)} e la \VUgras{schiumarola} \textsuperscript{(22)}.

%%K t06-tit-8 sub
\VUsaut{4.0mm}
\VUpk{10.2pt}{40}{La batteria da cucina.}
\VUsaut{3.0mm}

%%K t06-c4-06-1 p
6. --- La \VUgras{batteria di casseruole} \textsuperscript{(25)} appesa in fondo alla
stanza è molto pulita. Le belle \VUgras{casseruole} \textsuperscript{(26)} di rame, il
\VUgras{bricco del latte} \textsuperscript{(27)}, il \VUgras{colino}
\textsuperscript{(28)} e la \VUgras{grattugia} \textsuperscript{(29)} sono stati puliti
con cura.

%%K t06-c4-07-1 p
7. --- La \VUgras{saliera} \textsuperscript{(30)} sta sulla piccola credenza accanto
alla \VUgras{teiera} \textsuperscript{(31)} e alla \VUgras{damigiana dell'olio}
\textsuperscript{(32)}. Nel \VUgras{cassetto} \textsuperscript{(33)} ci saranno le
posate e i coltelli da cucina.

%%K t06-c4-08-1 p
8. --- La \VUgras{scopa} \textsuperscript{(34)} e il \VUgras{piumino}
\textsuperscript{(35)} sono in un angolo. La cuoca ha appena usato il \VUgras{ceppo da
macellaio} \textsuperscript{(36)}; l'ha lasciato vicino alla finestra.

%%K t06-c4-09-1 p
9. --- Un \VUgras{asciugamano} \textsuperscript{(37)} è appeso al muro accanto
all'\VUgras{imbuto} \textsuperscript{(38)}. In questa parte della cucina si tengono la
\VUgras{graticola} \textsuperscript{(39)}, la \VUgras{mannaia}
\textsuperscript{(40)} e il \VUgras{tagliere} \textsuperscript{(41)}. Poi, sopra la
mannaia, su una \VUgras{mensola} \textsuperscript{(42)}, ci sono \VUgras{pentole}
\textsuperscript{(43)}, la \VUgras{casseruola da stufato} \textsuperscript{(44)} con il
suo \VUgras{coperchio} \textsuperscript{(45)} e il \VUgras{paiolo}
\textsuperscript{(46)}. La \VUgras{padella} \textsuperscript{(47)} è appesa lì vicino.

%%K t06-c4-10-1 p
10. --- Solo il \VUgras{rubinetto} \textsuperscript{(48)} di rame getta un lampo in
questo angolo. L'\VUgras{acquaio} \textsuperscript{(49)}, su cui sta la grande
\VUgras{brocca} \textsuperscript{(50)}, dev'essere molto comodo con il suo armadietto,
dove si tengono la \VUgras{botticella dell'aceto} \textsuperscript{(51)} con il suo
\VUgras{rubinetto} \textsuperscript{(52)}, la \VUgras{cassetta dei rifiuti}
\textsuperscript{(53)} e i \VUgras{piatti sporchi} \textsuperscript{(54)}.

%%K t06-tit-9 sub
\VUsaut{4.0mm}
\VUpk{10.2pt}{40}{Che cos'altro si tiene in cucina.}
\VUsaut{3.0mm}

%%K t06-c4-11-1 p
11. --- Il grande \VUgras{mastello} \textsuperscript{(55)} in cui si lava la
\VUgras{verdura} \textsuperscript{(58)}, e il \VUgras{paiolo} \textsuperscript{(56)} di
ghisa posato sul \VUgras{pavimento di piastrelle} \textsuperscript{(57)}, saranno di
quell'armadietto anche loro.

%%K t06-c4-12-1 p
12. --- Di fuochi la cuoca non ne manca di certo, perché qui sull'angolo della tavola
c'è un \VUgras{fornello a gas} \textsuperscript{(59)} con dell'acqua che scalda in un
pentolino. Il \VUgras{vapore} \textsuperscript{(60)} che ne esce ci dice che deve
bollire. Quest'acqua sarà per il caffè, visto che la \VUgras{caffettiera}
\textsuperscript{(64)} e il \VUgras{macinino} \textsuperscript{(63)} stanno all'altro
capo della tavola.

%%K t06-c4-13-1 p
13. --- Il fornello a gas è collegato al \VUgras{becco del gas}
\textsuperscript{(62)} da un lungo \VUgras{tubo di gomma} \textsuperscript{(61)}.

%%K t06-c4-14-1 p
14. --- La \VUgras{donna a ore} \textsuperscript{(66)}, che viene ad aiutare la cuoca,
prepara la verdura per la cena. Quando sarà sbucciata e lavata, la metterà nella
\VUgras{casseruola} \textsuperscript{(65)} fino al momento di cuocerla.

%%K t06-tit-10 sub
\VUsaut{4.0mm}
{\centering\fontsize{10.2pt}{10.2pt}\selectfont\textls[40]{\scshape La stanza da bagno.} \textit{(Il locale del bagno.)}\par}
\VUsaut{3.0mm}

%%K t06-c5-01-1 p
1. --- Ancora un'ultima curiosità e conosceremo le stanze principali della casa: gli
impianti della stanza da bagno. Non ci vorrà molto, perché è piccola, e vediamo subito
che la \VUgras{vasca} \textsuperscript{(1)} ha un buon \VUgras{scaldabagno}
\textsuperscript{(2)}, che il \VUgras{portasapone} \textsuperscript{(3)} è a portata di
mano di chi fa il bagno, e che il \VUgras{portasciugamani} \textsuperscript{(5)}
dev'essere comodo per asciugare gli \VUgras{asciugamani} \textsuperscript{(4)}.

%%K t06-c5-02-1 p
2. --- Le pareti di questa stanza sono rivestite di \VUgras{piastrelle smaltate}
\textsuperscript{(6)}, così è stato possibile installare una \VUgras{doccia}
\textsuperscript{(7)} senza che l'acqua che schizza durante l'uso le rovini. Un piccolo
\VUgras{catino} \textsuperscript{(8)} di zinco per i pediluvi completa
l'arredamento.
\VUsaut{6.0mm}
\VUfilet{22mm}
